% ============================================================
%  C: Como Programar -- 6ª Edição | Deitel & Deitel
%  Capítulo 12 -- Estruturas de Dados em C
%  REFLEXÃO FINAL + SEÇÃO ESPECIAL DO COMPILADOR
%  (Cap. 12 não tem seção "Fazendo a Diferença" tradicional)
% ============================================================
\documentclass[12pt,a4paper]{article}
\usepackage[utf8]{inputenc}
\usepackage[T1]{fontenc}
\usepackage[brazil]{babel}
\usepackage{geometry}
\geometry{top=2.5cm,bottom=2.5cm,left=3cm,right=3cm}
\usepackage{parskip}\setlength{\parskip}{6pt}\setlength{\parindent}{0pt}
\usepackage{xcolor,tcolorbox,enumitem,listings,fancyhdr,titlesec,hyperref,booktabs,array,amsmath}
\tcbuselibrary{skins,breakable}

\definecolor{deitelblue}{RGB}{0,70,127}
\definecolor{deitelgray}{RGB}{240,242,245}
\definecolor{deitelgreen}{RGB}{0,110,60}
\definecolor{deitellightblue}{RGB}{220,235,250}
\definecolor{deitelred}{RGB}{160,20,20}
\definecolor{deitellorange}{RGB}{200,90,0}
\definecolor{ecogreen}{RGB}{0,120,50}
\definecolor{ecolightgreen}{RGB}{210,240,220}

\newtcolorbox{boaprática}[1][]{colback=deitellightblue,colframe=deitelblue,
  fonttitle=\bfseries\small,title={Boa prática de programação},breakable,#1}
\newtcolorbox{respostabox}[1][]{colback=deitelgray,colframe=deitelblue!60,
  fonttitle=\bfseries\small\color{deitelblue},title={Resposta detalhada},breakable,#1}
\newtcolorbox{enunciadoBox}[1][]{colback=white,colframe=deitelblue,
  fonttitle=\bfseries,title={#1},breakable}
\newtcolorbox{saida}[1][]{colback=black!88,colframe=deitelblue,
  fontupper=\ttfamily\small\color{green!70},
  title={\small\color{white}Saída do programa},breakable,#1}
\newtcolorbox{pesquisa}[1][]{colback=ecolightgreen,colframe=ecogreen,
  fonttitle=\bfseries\small,title={Pesquisa externa realizada},breakable,#1}
\newtcolorbox{atencao}[1][]{colback=yellow!10,colframe=deitellorange,
  fonttitle=\bfseries\small,title={Atenção},breakable,#1}

\setlist[enumerate,1]{label=\textbf{\alph*)},leftmargin=2em}
\setlist[itemize]{leftmargin=2em}

\lstset{
  basicstyle=\ttfamily\small,backgroundcolor=\color{deitelgray},
  frame=single,framerule=0.4pt,rulecolor=\color{deitelblue!40},
  breaklines=true,showstringspaces=false,
  keywordstyle=\color{deitelblue}\bfseries,
  commentstyle=\color{deitelgreen}\itshape,
  stringstyle=\color{deitelred},
  language=C,numbers=left,numberstyle=\tiny\color{gray},
  xleftmargin=18pt,xrightmargin=5pt,stepnumber=1,
}

\pagestyle{fancy}\fancyhf{}
\fancyhead[L]{\small\color{deitelblue}\textbf{C: Como Programar -- 6ª ed. | Deitel \& Deitel}}
\fancyhead[R]{\small\color{deitelblue}Cap. 12 -- Seção Especial}
\fancyfoot[C]{\small\thepage}
\renewcommand{\headrulewidth}{0.4pt}

\titleformat{\section}[block]{\large\bfseries\color{deitelblue}}{\thesection.}{0.5em}{}[\titlerule]
\titleformat{\subsection}[block]{\normalsize\bfseries\color{deitelblue!80}}{}{}{}

\hypersetup{colorlinks=true,linkcolor=deitelblue,urlcolor=ecogreen}

\begin{document}

\begin{titlepage}
  \centering\vspace*{2cm}
  {\color{deitelblue}\rule{\linewidth}{2pt}}\\[0.4cm]
  {\huge\bfseries\color{deitelblue}C: Como Programar}\\[0.2cm]
  {\large\color{deitelblue}6ª Edição $\cdot$ Paul Deitel \& Harvey Deitel}\\[0.2cm]
  {\color{deitelblue}\rule{\linewidth}{2pt}}\\[1cm]
  {\LARGE\bfseries Capítulo 12}\\[0.3cm]
  {\Large\color{deitelblue!80}Estruturas de Dados em C}\\[0.8cm]
  \begin{tcolorbox}[colback=ecolightgreen,colframe=ecogreen,width=0.85\linewidth]
    \centering{\large\bfseries\color{ecogreen}Seção Especial}\\[0.3cm]
    {\normalsize Construindo seu próprio compilador\\
    (Exercícios 12.26 -- 12.30)\\[0.2cm]
    \textbf{Reflexão final do capítulo}}
  \end{tcolorbox}
  \vfill{\small Abordagem incremental Deitel}
\end{titlepage}

\tableofcontents\newpage

% ============================================================
\section*{Nota Introdutória}

\begin{atencao}
O Capítulo~12 do livro \textbf{não tem seção ``Fazendo a Diferença''} tradicional.
Em seu lugar, encontra-se a \textbf{Seção Especial sobre Construção de Compilador}
(exercícios 12.26 a 12.30), que constitui um projeto de várias semanas integrando
\textit{todos} os conceitos vistos até aqui.

Este documento apresenta:
\begin{itemize}[noitemsep]
  \item Visão geral dos exercícios 12.26--12.30 (texto integral disponível em
  PDF separado no site da Deitel).
  \item Esqueletos de implementação dos componentes principais.
  \item Reflexão final consolidando todo o conhecimento do livro até o Cap.~12.
\end{itemize}
\end{atencao}

% ============================================================
\section{Visão Geral do Projeto do Compilador}
% ============================================================

\begin{enunciadoBox}[Contexto integrado]
O projeto integra três trabalhos anteriores:
\begin{itemize}[noitemsep]
  \item \textbf{Cap.~7}: simulador Simpletron e linguagem SML (exercícios 7.27--7.29).
  \item \textbf{Cap.~11}: leitura/escrita de programas em arquivos (Ex.~11.17).
  \item \textbf{Cap.~12}: estruturas de dados (pilhas, filas, listas, árvores) que
  o compilador precisa para processar código.
\end{itemize}

\textbf{O objetivo:} criar um compilador que traduza programas escritos em uma
linguagem de alto nível chamada \textbf{Simple} (variante didática de BASIC)
para Simpletron Machine Language (SML), que então executa no simulador do Cap.~7.
\end{enunciadoBox}

\subsection*{A linguagem Simple}

\begin{respostabox}
\textbf{Comandos da linguagem Simple} (sintaxe estilo BASIC):

\begin{lstlisting}
10 rem este e um comentario
20 input x          /* le inteiro do teclado e armazena em x */
30 let y = x * 2    /* atribuicao com expressao */
40 if y > 100 goto 60
50 print y
60 goto 20
70 end
\end{lstlisting}

\textbf{Características:}
\begin{itemize}[noitemsep]
  \item Cada linha começa com um \textbf{número de linha} (rótulo).
  \item Tipos: apenas inteiros (mapeados a palavras SML).
  \item Variáveis: identificadores simples (\texttt{a}, \texttt{x}, \texttt{i}, \ldots).
  \item Operadores: \texttt{+}, \texttt{-}, \texttt{*}, \texttt{/}, \texttt{\%},
  comparadores \texttt{<}, \texttt{>}, \texttt{<=}, \texttt{>=}, \texttt{==},
  \texttt{!=}.
  \item Controle: \texttt{goto}, \texttt{if-goto}.
\end{itemize}
\end{respostabox}

\subsection*{Exemplo de programa Simple}

\begin{respostabox}
Calcular a média de inteiros lidos (sentinela $= -1$):

\begin{lstlisting}
10  rem programa de media com sentinela
20  rem total acumulado em t, contador em c
30  let t = 0
40  let c = 0
50  input n
60  if n == -1 goto 110
70  let t = t + n
80  let c = c + 1
90  goto 50
100 rem calcula e imprime a media
110 if c == 0 goto 140
120 let m = t / c
130 print m
140 end
\end{lstlisting}
\end{respostabox}

\newpage

% ============================================================
\section{Exercício 12.26 -- Parser de Simple}
% ============================================================

\begin{respostabox}
\textbf{Etapas de um parser:}
\begin{enumerate}[noitemsep]
  \item \textbf{Análise léxica (\textit{tokenization})}: separar a entrada em
  tokens (palavras-chave, identificadores, números, operadores).
  \item \textbf{Análise sintática}: verificar se a sequência de tokens segue a
  gramática da linguagem.
  \item \textbf{Geração de tabela de símbolos}: mapear identificadores e rótulos
  para posições em memória.
\end{enumerate}

\textbf{Esqueleto da tabela de símbolos (\texttt{struct} do Cap.~10):}
\begin{lstlisting}
typedef struct {
    int  simbolo;     /* numero de linha ou ASCII de variavel */
    char tipo;        /* 'C' = constante, 'L' = linha, 'V' = variavel */
    int  posicao;     /* posicao em memoria SML (0-99) */
} Simbolo;

Simbolo tabela[ 100 ];
int tamanhoTabela = 0;
\end{lstlisting}

\textbf{Parser principal (esqueleto):}
\begin{lstlisting}
typedef struct linhaSimple {
    int   numeroLinha;
    char  comando[ 100 ];
    int   posicaoSML;          /* preenchido na primeira passada */
    struct linhaSimple *prox;
} LinhaSimple;

LinhaSimple* parseArquivo( const char *nomeArq )
{
    FILE *fp = fopen( nomeArq, "r" );
    LinhaSimple *lista = NULL, *ultimo = NULL;
    char buf[ 200 ];

    while ( fgets( buf, sizeof( buf ), fp ) != NULL ) {
        LinhaSimple *l = malloc( sizeof( LinhaSimple ) );
        sscanf( buf, "%d %[^\n]", &l->numeroLinha, l->comando );
        l->posicaoSML = -1;
        l->prox       = NULL;
        if ( lista == NULL ) lista = l;
        else                 ultimo->prox = l;
        ultimo = l;
    }
    fclose( fp );
    return lista;
}
\end{lstlisting}
\end{respostabox}

\newpage

% ============================================================
\section{Exercício 12.27 -- Geração de Código SML}
% ============================================================

\begin{respostabox}
\textbf{Estratégia de duas passadas:}

\textbf{Primeira passada:} percorrer todas as linhas Simple e:
\begin{itemize}[noitemsep]
  \item Atribuir uma posição SML para cada \textit{linha} (rótulo).
  \item Atribuir uma posição SML para cada \textit{variável} encontrada.
  \item \textbf{Não} resolver ainda os \texttt{goto} (eles podem apontar
  para frente).
\end{itemize}

\textbf{Segunda passada:} re-percorrer e gerar instruções SML reais, agora com
todos os endereços conhecidos.

\textbf{Mapeamento típico (Simple $\to$ SML):}

\begin{center}
\renewcommand{\arraystretch}{1.3}
\begin{tabular}{l l}
\toprule
\textbf{Simple} & \textbf{SML gerado} \\ \midrule
\texttt{rem ...}              & nada \\
\texttt{input x}              & \texttt{READ x} \\
\texttt{print x}              & \texttt{WRITE x} \\
\texttt{let x = expr}         & avalia \texttt{expr}, \texttt{STORE x} \\
\texttt{goto NL}              & \texttt{BRANCH posSML(NL)} \\
\texttt{if a == b goto NL}    & \texttt{LOAD a, SUBTRACT b, BRANCHZERO NL} \\
\texttt{if a < b goto NL}     & \texttt{LOAD a, SUBTRACT b, BRANCHNEG NL} \\
\texttt{end}                  & \texttt{HALT} \\
\bottomrule
\end{tabular}
\end{center}

\textbf{Tradução de expressões com pilha (Ex.~12.12!):}

A expressão \texttt{let z = (a + b) * 2} converte-se em pós-fixo
(\texttt{a b + 2 *}) e gera SML:
\begin{lstlisting}
LOAD a
ADD  b
STORE temp1
LOAD temp1
MULTIPLY const2  /* const2 = posicao da constante 2 */
STORE z
\end{lstlisting}

\textbf{Aqui está a beleza:} o Ex.~12.12 (conversão infixo $\to$ pós-fixo) e o
Ex.~12.13 (avaliação pós-fixa) que pareciam exercícios isolados são, na verdade,
\textbf{componentes essenciais} de um compilador real!
\end{respostabox}

\newpage

% ============================================================
\section{Exercícios 12.28 a 12.30 -- Extensões}
% ============================================================

\begin{respostabox}
\textbf{12.28 -- Melhoria do compilador:} otimizações comuns:
\begin{itemize}[noitemsep]
  \item \textbf{Constant folding:} avaliar expressões constantes em tempo de
  compilação (\texttt{let x = 2 + 3} $\to$ \texttt{let x = 5}).
  \item \textbf{Eliminação de código morto:} remover linhas inacessíveis.
  \item \textbf{Compartilhamento de constantes:} se a constante 2 aparece várias
  vezes, ela ocupa apenas uma posição SML.
\end{itemize}

\textbf{12.29 -- Recursos avançados:}
\begin{itemize}[noitemsep]
  \item Comandos compostos com chaves (\texttt{\{ ... \}}).
  \item Estruturas \texttt{while-end} e \texttt{for-end}.
  \item Tipos adicionais (ponto flutuante, strings curtas).
  \item Funções definidas pelo usuário (precisaria expandir SML também).
\end{itemize}

\textbf{12.30 -- Testes integrados:}
\begin{enumerate}[noitemsep]
  \item Escrever programas Simple.
  \item Compilá-los com o compilador.
  \item Executá-los no Simpletron.
  \item Verificar saídas esperadas.
\end{enumerate}

Este é o \textbf{ciclo completo} da computação: linguagem de alto nível $\to$
compilação $\to$ código de máquina $\to$ execução. Compreender esse ciclo é o
que separa um programador básico de um cientista da computação.
\end{respostabox}

\bigskip

% ============================================================
\section*{Reflexão Final -- O Encerramento do Núcleo de C}

\begin{tcolorbox}[colback=deitellightblue,colframe=deitelblue,
  title={\bfseries 12 capítulos, uma jornada completa}]

Chegando ao Capítulo~12, você completou \textbf{o núcleo essencial de C}.
Olhando para trás:

\textbf{Capítulos 1--2 -- Fundamentos:}
\begin{itemize}[noitemsep]
  \item Computadores, hardware, software.
  \item Primeiros programas: \texttt{printf}, \texttt{scanf}, \texttt{if}.
\end{itemize}

\textbf{Capítulos 3--5 -- Programação Estruturada:}
\begin{itemize}[noitemsep]
  \item \texttt{while}, \texttt{for}, \texttt{do-while}, \texttt{switch}.
  \item Operadores lógicos, atribuição composta.
  \item Funções, recursão, escopo, números aleatórios.
\end{itemize}

\textbf{Capítulos 6--8 -- Dados Compostos:}
\begin{itemize}[noitemsep]
  \item Arrays uni e bidimensionais.
  \item Ponteiros: o coração de C.
  \item Strings e manipulação de caracteres.
\end{itemize}

\textbf{Capítulo 9 -- E/S Refinada:}
\begin{itemize}[noitemsep]
  \item Domínio completo de \texttt{printf} e \texttt{scanf}.
\end{itemize}

\textbf{Capítulos 10--12 -- Estruturas e Persistência:}
\begin{itemize}[noitemsep]
  \item \texttt{struct}, \texttt{union}, manipulação de bits.
  \item Arquivos sequenciais e de acesso aleatório.
  \item Listas, pilhas, filas, árvores -- dinâmicas e poderosas.
\end{itemize}

\textbf{Você é agora capaz de:}
\begin{itemize}[noitemsep]
  \item Modelar problemas complexos com tipos próprios.
  \item Implementar estruturas de dados clássicas do zero.
  \item Persistir informações em arquivos.
  \item Construir programas significativos: jogos, simulações, calculadoras,
  analisadores de texto, e até compiladores.
\end{itemize}

\textbf{O que ainda falta?} Os capítulos 13--14 do livro abordam:
\begin{itemize}[noitemsep]
  \item \textbf{Cap.~13:} pré-processador (\texttt{\#define}, \texttt{\#ifdef},
  macros).
  \item \textbf{Cap.~14:} tópicos avançados (\texttt{va\_list}, manipulação de
  sinais, alocação dinâmica avançada, processamento de linha de comando).
\end{itemize}

Mas, em essência, \textbf{você já domina C}. Está pronto para:
\begin{itemize}[noitemsep]
  \item Ler código-fonte do Linux, do Git, do Vim.
  \item Contribuir com projetos open-source escritos em C.
  \item Aprender outras linguagens (C++, Rust, Go) com facilidade -- elas
  herdam grande parte da sintaxe e dos conceitos de C.
  \item Trabalhar com sistemas embarcados, drivers de dispositivos, kernels.
\end{itemize}

\end{tcolorbox}

\newpage

\begin{tcolorbox}[colback=ecolightgreen,colframe=ecogreen,
  title={\bfseries Estruturas de dados: o ponto de virada da carreira}]

Os conceitos do Capítulo~12 não são apenas mais um tópico -- são um
\textbf{ponto de virada} na sua maturidade técnica:

\textbf{Antes:} ``como faço para resolver este problema específico?''

\textbf{Depois:} ``qual estrutura de dados é mais adequada para esta família de
problemas?''

\textbf{Onde estruturas de dados aparecem na carreira:}
\begin{itemize}[noitemsep]
  \item \textbf{Entrevistas de emprego:} 80\% das perguntas em FAANG envolvem
  listas, árvores, hashes, grafos.
  \item \textbf{Code review:} ``por que você usou array aqui em vez de hash table?''
  \item \textbf{Performance:} a diferença entre $O(n)$ e $O(\log n)$ pode ser
  a diferença entre uma aplicação que funciona e outra que falha sob carga.
  \item \textbf{Design de sistemas:} bancos de dados usam B-trees; caches usam
  hash tables; redes usam filas; UNDO usa pilhas.
\end{itemize}

\textbf{Próximos níveis (autoestudo):}
\begin{itemize}[noitemsep]
  \item \textbf{Estruturas avançadas:} grafos, heaps, tries, segment trees.
  \item \textbf{Análise de algoritmos:} notação Big-O, master theorem, amortização.
  \item \textbf{Estruturas balanceadas:} AVL, Red-Black, B-trees.
  \item \textbf{Persistência:} estruturas funcionais imutáveis.
  \item \textbf{Concorrência:} estruturas lock-free, CAS, transactional memory.
\end{itemize}

\textbf{Livros recomendados:}
\begin{itemize}[noitemsep]
  \item Cormen et al., \textit{Introduction to Algorithms} (CLRS).
  \item Sedgewick, \textit{Algorithms in C}.
  \item Knuth, \textit{The Art of Computer Programming} (vol. 1, 3).
\end{itemize}
\end{tcolorbox}

\bigskip

\begin{boaprática}
\textbf{Conselho final do Cap.~12:}

A programação é, no fundo, sobre \textit{estruturar informação}. Tipos
primitivos (\texttt{int}, \texttt{double}) representam fatos isolados;
arrays representam coleções homogêneas; \texttt{struct} representa entidades
do mundo real; listas e árvores representam relações entre essas entidades.

Quando você domina essa hierarquia -- de bits a estruturas dinâmicas -- você
pode modelar \textit{qualquer coisa}. Os 12 capítulos da progressão Deitel
construíram essa hierarquia passo a passo, sem atalhos.

\textbf{O segredo da Programação} (com P maiúsculo): não é decorar sintaxe.
É treinar a mente para reconhecer padrões, dividir problemas, escolher
representações adequadas, e expressar soluções em código limpo e correto.

Você está pronto para o próximo nível. Bons estudos!
\end{boaprática}

\end{document}
